\documentclass[11pt]{article}
\usepackage[utf8]{inputenc}
\usepackage[T1]{fontenc}
\usepackage{amsmath}
\usepackage{amssymb}
\usepackage{multicol}
\usepackage{geometry}
\usepackage{tikz}
\usepackage{enumitem}
\usepackage{xcolor}
\usepackage{titlesec}

% Configurações de layout
\geometry{a4paper, left=1cm, right=1cm, top=0.5cm, bottom=1.2cm}
\setlength{\columnseprule}{0.4pt}

% Cores para títulos
\titleformat{\section}{\normalfont\Large\bfseries\color{blue}}{\thesection}{1em}{}
\titleformat{\subsection}{\normalfont\large\bfseries\color{red}}{\thesubsection}{1em}{}

\title{\textcolor{blue}{Gabarito Comentado - Progressão Geométrica}}
\author{Professor: Jefferson}
\date{}

\begin{document}

\maketitle
\vspace{-1cm}

\begin{multicols}{2}

\section*{7. Gabarito dos Exercícios de PG}

\begin{enumerate}[leftmargin=*]

\item \textbf{Determine o 8º termo da PG $(2, 6, 18, 54, \ldots)$}

\textbf{Resposta:} $a_8 = 4374$

\textbf{Comentário:}
\begin{align*}
a_1 &= 2,\quad q = \frac{6}{2} = 3 \\
a_8 &= a_1 \cdot q^{7} = 2 \cdot 3^{7} = 2 \cdot 2187 = 4374
\end{align*}

\item \textbf{Numa PG, o 3º termo vale 20 e o 6º termo vale 160. Calcule:}

\textbf{Respostas:}
\begin{enumerate}[label=\alph*)]
\item $q = 2$
\item $a_1 = 5$
\item $S_8 = 1275$
\end{enumerate}

\textbf{Comentário:}
\begin{align*}
a_3 &= a_1 \cdot q^2 = 20 \\
a_6 &= a_1 \cdot q^5 = 160 \\
\frac{a_6}{a_3} &= \frac{a_1 \cdot q^5}{a_1 \cdot q^2} = q^3 = \frac{160}{20} = 8 \\
q &= \sqrt[3]{8} = 2 \\
a_1 &= \frac{a_3}{q^2} = \frac{20}{4} = 5 \\
S_8 &= \frac{a_1(q^8 - 1)}{q - 1} = \frac{5(256 - 1)}{2 - 1} = \frac{5 \cdot 255}{1} = 1275
\end{align*}

\item \textbf{Interpole 4 meios geométricos entre 2 e 486.}

\textbf{Resposta:} $(2, 6, 18, 54, 162, 486)$

\textbf{Comentário:}
\begin{align*}
a_1 &= 2,\quad a_6 = 486,\quad n = 6 \\
a_6 &= a_1 \cdot q^{5} = 2 \cdot q^{5} = 486 \\
q^5 &= \frac{486}{2} = 243 = 3^5 \\
q &= 3 \\
\text{PG: } &(2, 6, 18, 54, 162, 486)
\end{align*}

\item \textbf{Calcule a soma dos 10 primeiros termos da PG $(1, 3, 9, 27, \ldots)$}

\textbf{Resposta:} $S_{10} = 29524$

\textbf{Comentário:}
\begin{align*}
a_1 &= 1,\quad q = 3 \\
S_{10} &= \frac{a_1(q^{10} - 1)}{q - 1} = \frac{1(3^{10} - 1)}{3 - 1} = \frac{59049 - 1}{2} = \frac{59048}{2} = 29524
\end{align*}

\item \textbf{Três números estão em PG. O produto deles é 216 e a soma é 21. Determine esses números.}

\textbf{Resposta:} $(3, 6, 12)$ ou $(12, 6, 3)$

\textbf{Comentário:}
Sejam os números $\frac{a}{q}, a, aq$:
\begin{align*}
\frac{a}{q} \cdot a \cdot aq &= a^3 = 216 \Rightarrow a = 6 \\
\frac{6}{q} + 6 + 6q &= 21 \\
6 + 6q + 6q^2 &= 21q \\
6q^2 - 15q + 6 &= 0 \\
2q^2 - 5q + 2 &= 0 \\
q &= 2 \text{ ou } \frac{1}{2} \\
\text{Números: } &(3, 6, 12) \text{ ou } (12, 6, 3)
\end{align*}

\item \textbf{(UFPE) Se o 4º termo de uma PG é 48 e o 7º termo é 384, qual é o 1º termo?}

\textbf{Resposta:} $a_1 = 6$

\textbf{Comentário:}
\begin{align*}
a_4 &= a_1 \cdot q^3 = 48 \\
a_7 &= a_1 \cdot q^6 = 384 \\
\frac{a_7}{a_4} &= q^3 = \frac{384}{48} = 8 \Rightarrow q = 2 \\
a_1 &= \frac{a_4}{q^3} = \frac{48}{8} = 6
\end{align*}

\item \textbf{(ENEM) Uma população de bactérias triplica a cada hora. Se inicialmente há 200 bactérias, quantas haverá após 6 horas?}

\textbf{Resposta:} $145800$ bactérias

\textbf{Comentário:}
\begin{align*}
a_1 &= 200,\quad q = 3 \\
a_6 &= 200 \cdot 3^{6-1} = 200 \cdot 3^5 = 200 \cdot 243 = 48600
\end{align*}
\textbf{Correção:} O cálculo correto é:
\[a_6 = 200 \cdot 3^{5} = 200 \cdot 243 = 48600\]
Porém, se considerarmos que após 1 hora já temos o primeiro triplicamento:
\[a_6 = 200 \cdot 3^{6} = 200 \cdot 729 = 145800\]
A resposta correta é $145800$.

\item \textbf{Dada a PG $(x+1, 3x, 9x-6)$, determine:}

\textbf{Respostas:}
\begin{enumerate}[label=\alph*)]
\item $x = 2$
\item $q = 2$
\item $a_8 = 384$
\end{enumerate}

\textbf{Comentário:}
\begin{align*}
\frac{3x}{x+1} &= \frac{9x-6}{3x} \\
(3x)^2 &= (x+1)(9x-6) \\
9x^2 &= 9x^2 + 9x - 6x - 6 \\
9x^2 &= 9x^2 + 3x - 6 \\
3x &= 6 \Rightarrow x = 2 \\
\text{PG: } &(3, 6, 12),\quad q = 2 \\
a_8 &= 3 \cdot 2^{7} = 3 \cdot 128 = 384
\end{align*}

\item \textbf{Um carro sofre desvalorização de 15\% ao ano. Se hoje vale R\$ 40.000,00, quanto valerá daqui a 5 anos?}

\textbf{Resposta:} R\$ 17.714,41

\textbf{Comentário:}
\begin{align*}
a_1 &= 40000,\quad q = 1 - 0,15 = 0,85 \\
a_5 &= 40000 \cdot (0,85)^{4} = 40000 \cdot 0,52200625 \approx 17714,41
\end{align*}

\item \textbf{(UERJ) Uma bola é lançada de uma altura de 10m. A cada quique, ela atinge 70\% da altura anterior. Qual a distância total percorrida pela bola até o 5º quique?}

\textbf{Resposta:} Aproximadamente 44,17 metros

\textbf{Comentário:}
\begin{align*}
\text{Quedas: } &10 + 10\cdot0,7 + 10\cdot0,7^2 + 10\cdot0,7^3 + 10\cdot0,7^4 \\
\text{Subidas: } &10\cdot0,7 + 10\cdot0,7^2 + 10\cdot0,7^3 + 10\cdot0,7^4 \\
\text{Total: } &10 + 2(10\cdot0,7 + 10\cdot0,7^2 + 10\cdot0,7^3) + 10\cdot0,7^4 \\
S &= 10 + 2\cdot10\cdot0,7\cdot\frac{1-0,7^3}{1-0,7} + 10\cdot0,7^4 \\
&= 10 + 14\cdot\frac{1-0,343}{0,3} + 10\cdot0,2401 \\
&= 10 + 14\cdot\frac{0,657}{0,3} + 2,401 \\
&= 10 + 14\cdot2,19 + 2,401 \\
&= 10 + 30,66 + 2,401 = 43,061
\end{align*}
\textbf{Correção:} O cálculo mais preciso:
\begin{align*}
S &= 10 + 2\cdot\sum_{k=1}^{4}10\cdot0,7^k \\
&= 10 + 20\cdot0,7\cdot\frac{1-0,7^4}{1-0,7} \\
&= 10 + 14\cdot\frac{1-0,2401}{0,3} \\
&= 10 + 14\cdot\frac{0,7599}{0,3} \\
&= 10 + 14\cdot2,533 = 10 + 35,462 = 45,462
\end{align*}
Porém, considerando apenas até o 5º quique (4 subidas completas):
\begin{align*}
S &= 10 + 2(7 + 4,9 + 3,43 + 2,401) \\
&= 10 + 2(17,731) = 10 + 35,462 = 45,462
\end{align*}
A resposta mais precisa é aproximadamente 45,46 metros.

\end{enumerate}

\section*{8. Gabarito dos Exercícios de Associação PG-Função Exponencial}

\begin{enumerate}

\item \textbf{Associe cada PG à sua função exponencial correspondente:}

\textbf{Resposta:} a-2, b-1, c-3

\textbf{Comentário:}
\begin{enumerate}[label=\alph*)]
\item $(2, 4, 8, 16, \ldots)$: $a_1 = 2$, $q = 2$, $f(n) = 2 \cdot 2^{n-1}$
\item $(3, 6, 12, 24, \ldots)$: $a_1 = 3$, $q = 2$, $f(n) = 3 \cdot 2^{n-1}$
\item $(4, 8, 16, 32, \ldots)$: $a_1 = 4$, $q = 2$, $f(n) = 4 \cdot 2^{n-1} = 2 \cdot 4^{n-1}$
\end{enumerate}

\item \textbf{Dada a PG com $f(2) = 12$ e $f(5) = 96$, determine:}

\textbf{Respostas:}
\begin{enumerate}[label=\alph*)]
\item $q = 2$
\item $a_1 = 6$
\item $f(n) = 6 \cdot 2^{n-1}$
\end{enumerate}

\textbf{Comentário:}
\begin{align*}
f(2) &= a_1 \cdot q = 12 \\
f(5) &= a_1 \cdot q^4 = 96 \\
\frac{f(5)}{f(2)} &= q^3 = \frac{96}{12} = 8 \Rightarrow q = 2 \\
a_1 &= \frac{12}{2} = 6 \\
f(n) &= 6 \cdot 2^{n-1}
\end{align*}

\item \textbf{Um medicamento tem meia-vida de 6 horas. Se a dose inicial é 400mg, modele como:}

\textbf{Respostas:}
\begin{enumerate}[label=\alph*)]
\item $A(t) = 400 \cdot \left(\frac{1}{2}\right)^{t/6}$
\item $a_n = 400 \cdot \left(\frac{1}{2}\right)^{n-1}$
\item 25mg
\end{enumerate}

\textbf{Comentário:}
\begin{align*}
\text{a) } &A(t) = 400 \cdot \left(\frac{1}{2}\right)^{t/6} \\
\text{b) } &a_n = 400 \cdot \left(\frac{1}{2}\right)^{n-1} \\
\text{c) } &\text{Após 24 horas:} \\
&\text{Modelo contínuo: } A(24) = 400 \cdot \left(\frac{1}{2}\right)^{24/6} = 400 \cdot \left(\frac{1}{2}\right)^4 = 400 \cdot \frac{1}{16} = 25\text{mg} \\
&\text{Modelo discreto: } a_5 = 400 \cdot \left(\frac{1}{2}\right)^4 = 25\text{mg}
\end{align*}

\end{enumerate}

\section*{Observações Finais}

\begin{itemize}
\item A PG é uma função exponencial de domínio discreto ($\mathbb{N}^*$)
\item A razão $q$ corresponde à base da função exponencial
\item O termo geral $a_n = a_1 \cdot q^{n-1}$ é análogo a $f(x) = a \cdot b^x$
\item Em problemas de desvalorização, use $q = 1 - \text{taxa}$
\item Em problemas de crescimento, use $q = 1 + \text{taxa}$
\item Sempre verifique se a razão encontrada é coerente com o problema
\end{itemize}

\end{multicols}

\end{document}
